% GENERATED FILE. Edit canonical lessons, then run npm run generate:books.
% canonical-source-hash: fnv1a64:9f5835e545dc55d3
% canonical-lessons: TA-C204-patikkattu, TA-C204-maruntukkatai, TA-C204-unavakam, TA-C204-palporul-ankati, TA-C204-tunikkatai

\chapter{Stairs, Pharmacy, Restaurant, Supermarket, Clothes shop}
\label{ch:ta-stairs-pharmacy-restaurant-supermarket-clothes-shop}
\hlchaptermodality{204}

\begin{chapteropening}
\textbf{By the end of this chapter:} \emph{I can say stairs, a pharmacy, a restaurant, a supermarket and a clothes shop.}

Complete the last lesson of chapter 204: I can say stairs, a pharmacy, a restaurant, a supermarket and a clothes shop.
\end{chapteropening}

\section[\texorpdfstring{\ta{படிக்கட்டு} (paṭikkaṭṭu)}{படிக்கட்டு (paṭikkaṭṭu)}]{\ta{படிக்கட்டு} (paṭikkaṭṭu) — stairs, steps}

\label{lesson:TA-C204-patikkattu}



\begin{quote}
Before the new one: say the Tamil for until then, then the Tamil for childhood.
\end{quote}

\subsection*{You'll want to know: \ta{படிக்கட்டு}}
\textbf{\ta{படிக்கட்டு}} — \emph{paṭikkaṭṭu} — \textquotedblleft{}stairs, steps\textquotedblright{}.

\begin{grammarlens}[title={What you've built}]
One more word for this chapter.
\end{grammarlens}

\subsection*{Your turn}
\begin{itemize}
\raggedright
  \item \emph{Say it:} \emph{paṭikkaṭṭu}
  \item \emph{Say it:} \emph{paṭikkaṭṭu}, once more
  \item \emph{Say it:} say \emph{paṭikkaṭṭu}
  \item \emph{Recall:} say the Tamil for until then, then the Tamil for childhood, then say \emph{paṭikkaṭṭu} again
\end{itemize}

\begin{tcolorbox}[breakable,colback=teal!4,colframe=teal!35!black,title={Before you move on}]
What does \ta{படிக்கட்டு} mean? (\textquotedblleft{}Stairs, steps\textquotedblright{}.) Say it once more.
\end{tcolorbox}
\section[\texorpdfstring{\ta{மருந்துக்கடை} (maruntukkaṭai)}{மருந்துக்கடை (maruntukkaṭai)}]{\ta{மருந்துக்கடை} (maruntukkaṭai) — a pharmacy}

\label{lesson:TA-C204-maruntukkatai}



\begin{quote}
Before the new one: say the Tamil for childhood, then the Tamil for stairs.
\end{quote}

\subsection*{You'll want to know: \ta{மருந்துக்கடை}}
\textbf{\ta{மருந்துக்கடை}} — \emph{maruntukkaṭai} — \textquotedblleft{}a pharmacy\textquotedblright{}.

\begin{grammarlens}[title={What you've built}]
One more word for this chapter.
\end{grammarlens}

\subsection*{Your turn}
\begin{itemize}
\raggedright
  \item \emph{Say it:} \emph{maruntukkaṭai}
  \item \emph{Say it:} \emph{maruntukkaṭai}, once more
  \item \emph{Say it:} say \emph{maruntukkaṭai}
  \item \emph{Recall:} say the Tamil for childhood, then the Tamil for stairs, then say \emph{maruntukkaṭai} again
\end{itemize}

\begin{tcolorbox}[breakable,colback=teal!4,colframe=teal!35!black,title={Before you move on}]
What does \ta{மருந்துக்கடை} mean? (\textquotedblleft{}A pharmacy\textquotedblright{}.) Say it once more.
\end{tcolorbox}
\section[\texorpdfstring{\ta{உணவகம்} (uṇavakam)}{உணவகம் (uṇavakam)}]{\ta{உணவகம்} (uṇavakam) — a restaurant}

\label{lesson:TA-C204-unavakam}



\begin{quote}
Before the new one: say the Tamil for stairs, then the Tamil for a pharmacy.
\end{quote}

\subsection*{You'll want to know: \ta{உணவகம்}}
\textbf{\ta{உணவகம்}} — \emph{uṇavakam} — \textquotedblleft{}a restaurant\textquotedblright{}.

\begin{grammarlens}[title={What you've built}]
One more word for this chapter.
\end{grammarlens}

\subsection*{Your turn}
\begin{itemize}
\raggedright
  \item \emph{Say it:} \emph{uṇavakam}
  \item \emph{Say it:} \emph{uṇavakam}, once more
  \item \emph{Say it:} say \emph{uṇavakam}
  \item \emph{Recall:} say the Tamil for stairs, then the Tamil for a pharmacy, then say \emph{uṇavakam} again
\end{itemize}

\begin{tcolorbox}[breakable,colback=teal!4,colframe=teal!35!black,title={Before you move on}]
What does \ta{உணவகம்} mean? (\textquotedblleft{}A restaurant\textquotedblright{}.) Say it once more.
\end{tcolorbox}
\section[\texorpdfstring{\ta{பல்பொருள்} \ta{அங்காடி} (palporuḷ aṅkāṭi)}{பல்பொருள் அங்காடி (palporuḷ aṅkāṭi)}]{\ta{பல்பொருள்} \ta{அங்காடி} (palporuḷ aṅkāṭi) — a supermarket}

\label{lesson:TA-C204-palporul-ankati}



\begin{quote}
Before the new one: say the Tamil for a pharmacy, then the Tamil for a restaurant.
\end{quote}

\subsection*{You'll want to know: \ta{பல்பொருள்} \ta{அங்காடி}}
\textbf{\ta{பல்பொருள்} \ta{அங்காடி}} — \emph{palporuḷ aṅkāṭi} — \textquotedblleft{}a supermarket\textquotedblright{}.

\begin{grammarlens}[title={What you've built}]
One more word for this chapter.
\end{grammarlens}

\subsection*{Your turn}
\begin{itemize}
\raggedright
  \item \emph{Say it:} \emph{palporuḷ aṅkāṭi}
  \item \emph{Say it:} \emph{palporuḷ aṅkāṭi}, once more
  \item \emph{Say it:} say \emph{palporuḷ aṅkāṭi}
  \item \emph{Recall:} say the Tamil for a pharmacy, then the Tamil for a restaurant, then say \emph{palporuḷ aṅkāṭi} again
\end{itemize}

\begin{tcolorbox}[breakable,colback=teal!4,colframe=teal!35!black,title={Before you move on}]
What does \ta{பல்பொருள்} \ta{அங்காடி} mean? (\textquotedblleft{}A supermarket\textquotedblright{}.) Say it once more.
\end{tcolorbox}
\section[\texorpdfstring{\ta{துணிக்கடை} (tuṇikkaṭai)}{துணிக்கடை (tuṇikkaṭai)}]{\ta{துணிக்கடை} (tuṇikkaṭai) — a clothes shop}

\label{lesson:TA-C204-tunikkatai}



\begin{quote}
Before the new one: say the Tamil for a restaurant, then the Tamil for a supermarket.
\end{quote}

\subsection*{You'll want to know: \ta{துணிக்கடை}}
\textbf{\ta{துணிக்கடை}} — \emph{tuṇikkaṭai} — \textquotedblleft{}a clothes shop\textquotedblright{}.

\begin{grammarlens}[title={What you've built}]
One more word for this chapter.
\end{grammarlens}

\subsection*{Your turn}
\begin{itemize}
\raggedright
  \item \emph{Say it:} \emph{tuṇikkaṭai}
  \item \emph{Say it:} \emph{tuṇikkaṭai}, once more
  \item \emph{Say it:} say \emph{tuṇikkaṭai}
  \item \emph{Recall:} say the Tamil for a restaurant, then the Tamil for a supermarket, then say \emph{tuṇikkaṭai} again
\end{itemize}

\begin{tcolorbox}[breakable,colback=teal!4,colframe=teal!35!black,title={Before you move on}]
What does \ta{துணிக்கடை} mean? (\textquotedblleft{}A clothes shop\textquotedblright{}.) Say it once more.
\end{tcolorbox}
